\subsection{Conditional Probability \& Bayes}

\begin{definition}[Conditional Probability]
The probability of $A$ given that $B$ has occurred is
\[
P(A \mid B) = \frac{P(A \cap B)}{P(B)}, \qquad P(B) > 0.
\]
Conditioning on $B$ restricts attention to the outcomes inside $B$ and rescales their probabilities to sum to 1, which is the role of the denominator.
\end{definition}

\begin{theorem}[Bayes' Theorem]
Bayes' theorem reverses the direction of conditioning, updating a hypothesis as evidence arrives:
\[
P(A \mid B) = \frac{P(B \mid A) \, P(A)}{P(B)}.
\]
The four pieces are worth memorizing by name:
\begin{itemize}[itemsep=-1ex, label=--]
  \item $P(A \mid B)$ is the posterior, the updated belief in $A$ after seeing $B$
  \item $P(B \mid A)$ is the likelihood, how well $A$ explains the evidence
  \item $P(A)$ is the prior, the belief in $A$ before the evidence
  \item $P(B)$ is the marginal, the total probability of the evidence
\end{itemize}
\end{theorem}

\begin{keybox}[Odds form skips the marginal]
Dividing Bayes' theorem for two competing hypotheses cancels the shared marginal $P(B)$:
\[
\underbrace{\frac{P(A \mid B)}{P(A^c \mid B)}}_{\text{posterior odds}}
= \underbrace{\frac{P(B \mid A)}{P(B \mid A^c)}}_{\text{likelihood ratio}}
\times \underbrace{\frac{P(A)}{P(A^c)}}_{\text{prior odds}}.
\]
Each piece of evidence multiplies the prior odds by its likelihood ratio, so repeated independent evidence just multiplies again. This avoids computing $P(B)$ and is the fastest route through two-hypothesis update problems.
\end{keybox}

\medskip

\begin{theorem}[Law of Total Probability]
If $B_1, B_2, \dots , B_n$ partition the sample space (mutually exclusive and exhaustive), then
\[
P(A) = \sum_i P(A \mid B_i) \, P(B_i).
\]
\end{theorem}

This is how the marginal in Bayes' theorem is usually assembled. When $P(A)$ is not known directly but the conditional probabilities on each case are, summing those conditionals weighted by the case probabilities builds it. The two-way form $P(A) = P(A \mid C)P(C) + P(A \mid C^c)P(C^c)$ covers most interview uses.

\medskip

\begin{example}[Monty Hall Problem]
Assume the contestant chooses Door 1 and Monty opens Door 2 to reveal a goat. Let $A_i$ be the event the car is behind Door $i$ (prior $P(A_i) = \frac{1}{3}$), and $B$ be the event Monty opens Door 2.

The likelihoods $P(B|A_i)$ are:
\begin{itemize}[itemsep=-1ex, label=--]
    \item $P(B|A_1) = \frac{1}{2}$ (Monty randomly picks Door 2 or 3)
    \item $P(B|A_2) = 0$ (Monty cannot reveal the car)
    \item $P(B|A_3) = 1$ (Monty is forced to open Door 2)
\end{itemize}

The marginal probability of $B$ is:
\[
P(B) = \sum_{i=1}^3 P(B|A_i)P(A_i) = \left(\frac{1}{2}\right)\left(\frac{1}{3}\right) + (0)\left(\frac{1}{3}\right) + (1)\left(\frac{1}{3}\right) = \frac{1}{2}
\]

By Bayes' Theorem, the posterior probabilities of winning by staying ($A_1$) or switching ($A_3$) are:
\[
P(A_1|B) = \frac{P(B|A_1)P(A_1)}{P(B)} = \frac{\left(\frac{1}{2}\right)\left(\frac{1}{3}\right)}{\frac{1}{2}} = \frac{1/6}{1/2} = \frac{1}{3}
\]
\[
P(A_3|B) = \frac{P(B|A_3)P(A_3)}{P(B)} = \frac{(1)\left(\frac{1}{3}\right)}{\frac{1}{2}} = \frac{1/3}{1/2} = \frac{2}{3}
\]

Switching doubles the probability of winning.
\end{example}
